\section{A genuine finite phase on the nonzero coefficient-stratum remainder}
\label{sec:nc-coefficient-nonzero-remainder-phase}

For a smooth closed marked carrier $L.N$, let $F$ be the union of its actual
connected smooth components contained in $L.\mathrm{cosupp}$, and let
$R=L.N\setminus F$. The remainder datum $A$ keeps the SAME ideal and marking
and the SAME boundary, and has carrier $R$. Here the original boundary of
$L$ is empty. The construction uses the existing full coefficient-restriction
theorem, including its reverse implication on entire actual test histories.
It does not assert that the merely normal-crossings old carrier has become
SNC, or identify the new auxiliary canonical invariant with an old ordered one.

\begin{theorem}[The actual component remainder supplies its own domain and original open equality]
\label{hir:nc-coefficient-remainder-original-domain}
\lean{Hironaka.HironakaNormalCrossingsCoefficientRemainderPhase.contactRemainder_original_domain}
\leanok
For any closed smooth $L.N$, coherent $L.\mathcal I$, positive marking and empty
original boundary, $A$ is in the actual ambient-dimensional invariant domain,
$A.\mathrm{cosupp}=L.\mathrm{cosupp}\cap R$, and $R$ lies in the actual open
$G=X\setminus F$. The ENTIRE original records of $A$ and $L$ restrict equally
to $G$.
\uses{hir:original-all-level-contact-remainder-nonzero-geometry,hir:smooth-selected-components,hir:dimtwo-unselected-smooth-components}
\end{theorem}
\begin{proof}
The actual selected and remaining smooth component unions are disjoint closed
smooth sets and have union $L.N$. Their defining difference gives
$R\subseteq G$. Near a point of $R$, membership in $R$ and in $L.N$ agrees;
relative order therefore gives the asserted entire cosupport equality.
The genuine component-remainder theorem supplies coherence and nonvanishing
along every point of $R$, without a nonvanishing assumption on $L$. Its dimension
bound is at most the ambient dimension: at any smooth carrier point the
inclusion into the whole ambient manifold compares its regular dimension
with the ambient one. The empty boundary is SNC and transverse to $R$.
These prove every invariant-domain field. On $G$ the equality
$R=L.N\setminus F$ identifies the two entire restricted carriers; the ideal,
marking and empty ordered boundary were retained literally. This proves the
whole original marked-record equality.
\end{proof}

\begin{theorem}[An internally chosen canonical remainder resolution is a genuine parent phase]
\label{hir:nc-coefficient-remainder-finite-parent-phase}
\lean{Hironaka.HironakaNormalCrossingsCoefficientRemainderPhase.exists_canonical_remainder_parent_phase}
\leanok
Under the preceding natural hypotheses and compactness of
$L.\mathrm{cosupp}$ there is an actual finite tower $T$, canonical and terminal
for $A$, which is a genuine test tower of the SAME parent $L$. All its actual
centres lie in the literal strict transforms of $R$, and
\[
 L_T.\mathrm{cosupp}\cap R_T=\varnothing.
\]
No incoming test history, transformed-ideal identity, maximum, or terminal
certificate is required.
\uses{hir:nc-coefficient-remainder-original-domain,hir:principalization-compact-support,hir:original-open-carrier-test-history-extension,hir:original-open-equal-current-test-records}
\end{theorem}
\begin{proof}
The actual closed remainder makes $L.\mathrm{cosupp}\cap R$ compact. Apply the
proved compact-cosupport canonical driver to the actual domain $A$ from the
preceding theorem; it constructs a finite canonical terminal tower. Since
$A.N\subseteq G$, the actual open-carrier extension theorem constructs parent
testness on the SAME tower at every chronological step. Every actual centre
lies in the current $A$-carrier, hence in $R_T$ at its prefix.

The whole current records of $A$ and $L$ restrict equally to the actual
inverse-image open $G_T$. A point of $R_T$ projects into $R\subseteq G$, so is
in $G_T$. If it belonged to $L_T.\mathrm{cosupp}$, coherent cosupport restriction
and this entire record equality would place it in $A_T.\mathrm{cosupp}$.
The latter is empty by actual terminality. This proves the whole asserted
intersection is empty, including points on accumulated exceptional divisors.
Every history used here was internally constructed from the original data.
\end{proof}

\begin{theorem}[Original marked NC data produce a finite coefficient-stratum remainder phase]
\label{hir:nc-original-coefficient-remainder-finite-phase}
\lean{Hironaka.HironakaNormalCrossingsCoefficientRemainderPhase.exists_original_nc_coefficient_remainder_phase}
\leanok
Let $\mathcal I$ be any coherent ideal with positive mark $d$ and compact
original cosupport on a smooth $n$-fold. Let $H$ be an actual normal-crossings
carrier with branch multiplicity everywhere at most $r$, and put
$Z=\{x:\operatorname{mult}_H(x)\ge r\}$. For the SAME original datum
$M=(X,\varnothing,\mathcal I,d)$, set $L=\mathcal C(M)|_Z$, and take its actual
component remainder $R$ as above. There is an actual finite tower, internally
canonical and terminal for its remainder datum, which is simultaneously a
genuine test tower of $L$ and of the ORIGINAL $M$. Every centre lies in the
strict transforms of $R$ and of $Z$, and the SAME literal controlled transforms
satisfy
\[
 M_T.\mathrm{cosupp}\cap R_T=\varnothing.
\]
\uses{hir:nc-coefficient-remainder-finite-parent-phase,mc:coefficient-restriction-full,mc:coefficient-restriction-support,hir:general-nc-highest-multiplicity-smooth-centre}
\end{theorem}
\begin{proof}
The actual highest-stratum theorem makes $Z$ closed smooth. The coefficient
ideal is coherent, its marking is the product of the positive derivative
markings $d-i$ for $0\le i<d$, and its boundary is the original empty list.
The exact coefficient-restriction support theorem gives
$L.\mathrm{cosupp}=M.\mathrm{cosupp}\cap Z$, so it is compact. Apply the preceding
internally constructed parent phase to this literal $L$.

The full coefficient-restriction tower equivalence turns this actual $L$-test
history into an actual $M$-test history with every centre in the literal strict
transforms of $Z$. It also identifies its ENTIRE final coefficient cosupport
with $M_T.\mathrm{cosupp}\cap Z_T$. The genuine inclusion $R\subseteq Z$ persists
under strict transform. Thus a point of $M_T.\mathrm{cosupp}\cap R_T$ would belong
to $L_T.\mathrm{cosupp}\cap R_T$, already proved empty. All markings, germs,
centres and projections here are those of the SAME constructed histories.
This produces the nonzero-remainder phase; it makes no assertion that the
initially selected zero-along components have been treated or that all old
NC branches and born exceptional members already form an admissible algorithm.
\end{proof}
